\section{Organización de productos de IA: Taste y trabajo en dos etapas}

Cuando el producto pasa del diseño de flujos al diseño de contexto, la forma de organizar el trabajo también tiene que cambiar. Zhang Yueguang (张月光) dice que los límites de las decisiones se vuelven difusos, que Taste adquiere importancia y que el trabajo deja de ser lineal para organizarse en dos etapas. Aquí Taste no es un gusto estético subjetivo, sino la capacidad de juzgar si la salida del modelo es buena o mala, dónde están los límites de la experiencia de usuario y si una dirección de producto merece la apuesta.

\begin{figure}[H]
\centering
\includegraphics[width=0.96\textwidth]{ai-product-org-two-stage.png}
\caption{Organización de productos de IA: el trabajo lineal se convierte en trabajo en dos etapas.}
\end{figure}

\begin{knowledgebox}{Lectura de la figura: primero validar la capacidad, después empaquetar el producto}
La primera etapa de la figura es la viabilidad del modelo: context, prompt, herramientas, combinación de modelos, evaluación y juicio de taste; solo la segunda etapa es el empaquetado del producto: flujos, UI, modelo de negocio y lugar en la mente del usuario. El flujo de trabajo lineal tradicional coloca la definición de requisitos al principio, mientras que un equipo AI Native primero debe confirmar que la capacidad del modelo existe realmente.
\end{knowledgebox}

\subsection{Por qué importa Taste}

La salida de la IA es incierta, y la entrada del usuario también lo es. El gerente de producto no puede limitarse a decir «aquí va un botón»; tiene que juzgar si la salida es lo bastante buena, si es controlable, si hay usuarios dispuestos a seguir usándola y si puede convertirse en una puerta de un solo sentido. Este juicio a menudo no puede sustituirse por completo con métricas, porque un producto en su etapa temprana todavía no tiene métricas estables.

\begin{knowledgebox}{Términos clave: organización de productos de IA}
\begin{tabularx}{\textwidth}{p{0.24\textwidth}p{0.32\textwidth}X}
\toprule
Término & Explicación breve & Papel en este episodio \\
\midrule
Taste & Capacidad de juzgar la calidad de la salida, la dirección del producto y la experiencia de usuario & Capacidad central cuando los límites de las decisiones son difusos. \\
Context Design & Diseñar la información, las herramientas y las restricciones visibles para el modelo & Capa intermedia de los productos AI Native. \\
Prompt Practice & Método para construir la entrada y la descripción de la tarea & Es solo una parte de context design. \\
Two-stage Work & Primero validar la capacidad del modelo, después empaquetar el flujo del producto & Sustituye la colaboración lineal tradicional. \\
Evaluation & Juzgar si la salida es estable, utilizable e iterable & Conecta taste con la ingeniería. \\
\bottomrule
\end{tabularx}
\end{knowledgebox}

\subsection{Resumen de la sección}

Lo esencial en la organización de productos de IA es que producto, diseño, ingeniería y capacidad del modelo trabajen dentro del mismo ciclo de retroalimentación. Sin taste, el equipo se deja deslumbrar fácilmente por una demo; sin validación de ingeniería, taste se vuelve pura especulación.
